\section{Implementierung des Patches}
\label{sec:implementierung}

\subsection{Anpassung von \texttt{get()} und \texttt{get\_or\_404()}}
\label{sec:patch-base}

Die zentrale Änderung betrifft ausschliesslich die Datei
\texttt{app/shared/services/base.py}.
Beide Methoden erhalten dieselben zwei optionalen Keyword-Only-Parameter, die
bereits in \texttt{get\_multi()} und \texttt{count()} vorhanden sind:

\begin{itemize}
    \item \texttt{organization\_id: str | None = None}: wird nur ausgewertet, wenn
          das Modell das Attribut besitzt (\texttt{hasattr}-Guard, identisch zu
          \texttt{get\_multi()});
    \item \texttt{include\_deleted: bool = False}: filtert soft-gelöschte
          Datensätze, sofern das Modell \texttt{deleted\_at} kennt.
\end{itemize}

\texttt{get\_or\_404()} leitet beide Parameter an \texttt{get()} weiter und
liefert \texttt{404} zurück, wenn kein Treffer gefunden wird, also sowohl bei
nicht vorhandenen als auch bei organisationsfremden und soft-gelöschten
Datensätzen.

\texttt{delete()} wurde ebenfalls angepasst: es ruft intern \texttt{get()} auf
und gibt \texttt{organization\_id} weiter, damit auch DELETE-Operationen auf
fremde Objekte ins Leere laufen.

Der vollständige Diff der drei Methoden ist nachfolgend dargestellt:

\begin{lstlisting}[style=custompython, caption={base.py -- gepatchte Methoden get(), get\_or\_404() und delete()}]
# VORHER
def get(self, db: Session, id: str) -> ModelType | None:
    result = db.execute(select(self.model).where(self.model.id == id))
    return result.scalar_one_or_none()

def get_or_404(self, db: Session, id: str) -> ModelType:
    obj = self.get(db, id)
    if obj is None:
        raise NotFoundError(self.model.__name__, id)
    return obj

# NACHHER
def get(self, db, id, *, organization_id=None, include_deleted=False):
    query = select(self.model).where(self.model.id == id)
    if organization_id and hasattr(self.model, "organization_id"):
        query = query.where(self.model.organization_id == organization_id)
    if not include_deleted and hasattr(self.model, "deleted_at"):
        query = query.where(self.model.deleted_at.is_(None))
    return db.execute(query).scalar_one_or_none()

def get_or_404(self, db, id, *, organization_id=None, include_deleted=False):
    obj = self.get(db, id,
                   organization_id=organization_id,
                   include_deleted=include_deleted)
    if obj is None:
        raise NotFoundError(self.model.__name__, id)
    return obj

def delete(self, db, id, soft_delete=True, *, organization_id=None):
    obj = self.get(db, id, organization_id=organization_id)
    if obj is None:
        return None
    ...
\end{lstlisting}

Die neuen Parameter sind \emph{keyword-only} (nach \texttt{*})
und haben Standardwerte, sodass bestehende Aufrufer sich unverändert verhalten.
Nur Aufrufer, die
\texttt{organization\_id} übergeben (Schritt 2), bekommen den Schutz aktiviert.

\subsection{Anpassung aller Aufrufer}
\label{sec:patch-callers}

Alle 83 identifizierten \texttt{get\_or\_404()}-Aufrufer sowie die
betroffenen \texttt{delete()}-Aufrufe werden nach einem einheitlichen Muster angepasst:

\begin{lstlisting}[style=custompython, caption={Aufruf-Muster vor und nach dem Patch (Beispiel: Fahrzeug-Endpunkt)}]
# VORHER -- kein Mandanten-Check
vehicle = vehicle_service.get_or_404(db, vehicle_id)

# NACHHER -- organisation_id wird mitgegeben
vehicle = vehicle_service.get_or_404(
    db, vehicle_id, organization_id=current_user.organization_id
)
\end{lstlisting}

Das gleiche Muster gilt für \texttt{delete()}-Aufrufe sowie für interne
Service-Aufrufe, die ihrerseits \texttt{get\_or\_404()} verwenden.

\textbf{Betroffene Dateien der API-Schicht (24 Router-Dateien)}

\begin{itemize}
    \item \texttt{fleet/api/}: \texttt{vehicles.py} (7 Aufrufe),
          \texttt{tickets.py} (5), \texttt{maintenance.py} (9),
          \texttt{gps.py} (8), \texttt{fuel.py} (3),
          \texttt{fuel\_pumps.py} (4), \texttt{fuel\_pump\_deliveries.py} (3)
    \item \texttt{tools/api/}: \texttt{tools.py} (4),
          \texttt{tool\_assignments.py} (1), \texttt{tool\_locations.py} (3),
          \texttt{tool\_calibrations.py} (2), \texttt{tool\_categories.py} (3),
          \texttt{tool\_cases.py} (5), \texttt{consumables.py} (4)
    \item \texttt{sat/api/}: \texttt{reports.py} (3), \texttt{machines.py} (3),
          \texttt{customers.py} (3), \texttt{contacts.py} (2),
          \texttt{assistances.py} (3), \texttt{attachments.py} (1, DELETE)
    \item \texttt{shared/api/}: \texttt{employees.py} (4),
          \texttt{cost\_centers.py} (11), \texttt{documents.py} (5),
          \texttt{users.py} (3, GET/PATCH/DELETE \texttt{/\{user\_id\}})
\end{itemize}

\textbf{Sonderfall: \texttt{shared/api/notifications.py}}\\%
Das \texttt{Notification}-Modell besitzt kein \texttt{organization\_id}-Feld,
sondern ist über \texttt{user\_id} an den eingeloggten Nutzer gebunden.
Der Schutz erfolgt hier nicht über den \texttt{hasattr}-Guard in \texttt{get()},
sondern über einen expliziten Eigentümercheck nach dem Laden der Ressource
(\texttt{notification.user\_id != current\_user.id} → HTTP\,403).

\textbf{Ausnahmen (kein Fix nötig)}
\begin{itemize}
    \item \texttt{shared/api/organizations.py}: Admin-Only-Endpunkte ohne
          Mandanten-Kontext (globale Ressource).
\end{itemize}

\textbf{Service-interne Aufrufe der Service-Schicht (4 Methoden)}

Einige Service-Methoden rufen intern \texttt{get\_or\_404()} einer anderen
Service-Instanz auf. Diese wurden ebenfalls gepatcht:

\begin{lstlisting}[style=custompython, caption={Service-interner Cross-Service-Lookup mit org-Check}]
# fleet/services/fuel_pump_delivery.py -- create()
pump = fuel_pump_service.get_or_404(
    db, obj_in.pump_id, organization_id=obj_in.organization_id
)

# fleet/services/fuel.py -- create()
vehicle = vehicle_service.get_or_404(
    db, obj_in.vehicle_id, organization_id=obj_in.organization_id
)
\end{lstlisting}

Die Methoden \texttt{dispense\_fuel()}, \texttt{receive\_delivery()},
\texttt{adjust\_level()} und \texttt{update\_maintenance()} der Klasse
\texttt{FuelPumpService} sowie \texttt{return\_tool()} der Klasse
\texttt{ToolAssignmentService} erhielten einen optionalen Parameter
\texttt{organization\_id}, der an das interne \texttt{get\_or\_404()}
weitergeleitet wird. Die zugehörigen API-Endpunkte übergeben
\texttt{current\_user.organization\_id}.

\subsection{Behandlung von Modellen ohne \texttt{organization\_id}}
\label{sec:patch-no-org}

Zwei Modelle tragen kein \texttt{organization\_id}-Feld und werden
daher vom \texttt{hasattr}-Guard in \texttt{get()} nicht mit dem
Mandanten-Filter belegt:

\begin{itemize}
    \item \texttt{Notification}: nutzer-gebunden über \texttt{user\_id};
          Schutz erfolgt via Eigentümercheck in der Router-Schicht.
    \item \texttt{Organization}: Die Ressource \emph{ist} der Mandant selbst;
          Endpunkte sind ausschliesslich für Super-Admins zugänglich.
\end{itemize}

Das Modell \texttt{User} besitzt hingegen \texttt{organization\_id} und
profitiert vollständig vom \texttt{hasattr}-Guard; die zugehörigen
Endpunkte in \texttt{users.py} wurden ebenfalls mit
\texttt{organization\_id=current\_user.organization\_id} gesichert.

Der \texttt{hasattr}-Guard in \texttt{get()} stellt sicher, dass ein
versehentlicher Aufruf mit \texttt{organization\_id} bei Modellen ohne
dieses Feld vollständig ignoriert wird; es gibt keine Laufzeit-Ausnahme.
